%%%%%%%%%%%%%%%%%%%%%%%%%%%%%%%%%%%%%%%%
% Beamer Presentation
% LaTeX Template
% Version 1.0 (10/11/12)
%
% This template has been downloaded from:
% http://www.LaTeXTemplates.com
%
% License:
% CC BY-NC-SA 3.0 (http://creativecommons.org/licenses/by-nc-sa/3.0/)
%
%%%%%%%%%%%%%%%%%%%%%%%%%%%%%%%%%%%%%%%%%

%----------------------------------------------------------------------------------------
%	PACKAGES AND THEMES
%----------------------------------------------------------------------------------------

\documentclass[xcolor=dvipsnames, aspectratio=169]{beamer}

\mode<presentation> {

% The Beamer class comes with a number of default slide themes
% which change the colors and layouts of slides. Below this is a list
% of all the themes, uncomment each in turn to see what they look like.

\usetheme{Madrid} %Hannover
% As well as themes, the Beamer class has a number of color themes
% for any slide theme. Uncomment each of these in turn to see how it
% changes the colors of your current slide theme.
\useoutertheme{infolines} % Alternatively: miniframes, infolines, split
\useinnertheme{circles}
\definecolor{UBCblue}{rgb}{0.04706, 0.13725, 0.26667} % UBC Blue (primary)
\usecolortheme[named=UBCblue]{structure}
}

\usepackage{graphicx} % Allows including images
\usepackage{booktabs} % Allows the use of \toprule, \midrule and \bottomrule in tables
\usepackage{textpos}
\usepackage{caption}
\usepackage[utf8]{inputenc}
\usepackage[brazilian]{babel}
\usepackage{csquotes}
\usepackage{listings}
\setbeamertemplate{caption}[numbered]
\usepackage[style=abnt]{biblatex}
\addbibresource{bibliography.bib}
\PassOptionsToPackage{useregional}{datetime2}
\usepackage{xcolor}
\usepackage{amsmath}


\definecolor{codegreen}{rgb}{0,0.6,0}
\definecolor{codegray}{rgb}{0.5,0.5,0.5}
\definecolor{codepurple}{rgb}{0.58,0,0.82}
\definecolor{backcolour}{rgb}{0.95,0.95,0.92}
\definecolor{string-color}{rgb}{0.3333, 0.5254, 0.345}

\lstdefinestyle{mystyle}{
    backgroundcolor=\color{backcolour},   
    commentstyle=\color{codegreen},
    keywordstyle=\color{string-color},
    keywordstyle=[2]{\color{codepurple}},
    keywordstyle=[3]{\color{magenta}},
    numberstyle=\tiny\color{codegray},
    stringstyle=\color{codepurple},
    basicstyle=\ttfamily\tiny,
    breakatwhitespace=false,         
    breaklines=true,                 
    captionpos=b,                    
    keepspaces=true,                 
    numbers=left,                    
    numbersep=5pt,                  
    showspaces=false,                
    showstringspaces=false,
    showtabs=false,                  
    tabsize=2,
    % language=Python, % Adiciona suporte para Python
    morekeywords={If, Then, Else, While, Do, For, Return, end, End, if}, % Adicione suas palavras-chave aqui
    morecomment=[l]{//}, % Define o estilo de comentário, ajuste conforme necessário
    morecomment=[s]{/*}{*/}, % Para comentários de bloco, se necessário
    morestring=[b]" % Para strings, se necessário
}

\lstset{style=mystyle}
\newcommand{\source}[1]{\vspace{-20pt} \caption*{ Fonte: {#1}} }
\usepackage{copyrightbox}


\makeatletter
% \beamer@nav@subsectionstyle{hide/hide/hide}
\addtobeamertemplate{sidebar left}{%
\hspace{0.5cm}\includegraphics[width=0.9cm, keepaspectratio]{figures/_brasao_ufsm_cor.png}
% \hspace{2.3cm}\includegraphics[width=0.8cm, keepaspectratio]{figures/brasao_ctism.png}
% \hspace{2.3cm}\includegraphics[width=1.5cm, keepaspectratio]{_logosbc.png}
% \hspace{2.3cm}\includegraphics[width=1.5cm, keepaspectratio]{_logoERRC.png}%
}{}


\setbeamertemplate{footline}
{
	\leavevmode%
	\hbox{%
	    % \hspace{0.5cm}\includegraphics[width=0.8cm, keepaspectratio]{figures/_brasao_ufsm_cor.png}
		\begin{beamercolorbox}[wd=.333333\paperwidth,ht=2.25ex,dp=1ex,right]{date in head/foot}%
			\usebeamerfont{date in head/foot}\insertshortdate{}\hspace*{2em}
			\insertframenumber{} / \inserttotalframenumber\hspace*{2ex} 
		\end{beamercolorbox}}%
		%\vskip0pt%
	}
\makeatother

%----------------------------------------------------------------------------------------
%	TITLE PAGE
%----------------------------------------------------------------------------------------

\title[An experimental system for incremental environment modelling by an autonomous mobile robot]{An experimental system for incremental environment modelling by an autonomous mobile robot} % The short title appears at the bottom of every slide, the full title is only on the title page

\author[FDR]{Fábio Demo da Rosa} % Your name
%\includegraphics[]{logositeredes.png}
\institute[UFSM] % Your institution as it will appear on the bottom of every slide, may be shorthand to save space
{
Universidade Federal de Santa Maria \\ % Your institution for the title page
Pós-Graduação em Ciência da Computação \\
Disciplina de Robótica Móvel\\
\medskip
\textit{faberdemo@gmail.com} % Your email address
}
\date{\today} % Date, can be changed to a custom date
\newcounter{saveenumi}
\newcommand{\seti}{\setcounter{saveenumi}{\value{enumi}}}
\newcommand{\conti}{\setcounter{enumi}{\value{saveenumi}}}

\resetcounteronoverlays{saveenumi}


\begin{document}

\begin{frame}
\titlepage % Print the title page as the first slide
\end{frame}

\begin{frame}
\frametitle{Visão Geral} %\includegraphics[]{logositeredes.png}} % Table of contents slide, comment this block out to remove it
\tableofcontents % Throughout your presentation, if you choose to use \section{} and \subsection{} commands, these will automatically be printed on this slide as an overview of your presentation
\end{frame}

%----------------------------------------------------------------------------------------
%	PRESENTATION SLIDES
%----------------------------------------------------------------------------------------

%------------------------------------------------
\section{Introdução}
%------------------------------------------------
\begin{frame}[fragile]
  \frametitle{Introdução}
  \begin{itemize}
    \item Mapeamento incremental por robôs autônomos.
    \item Desafio de inconsistências e imprecisões nos dados dos sensores.
    \item Solução: Processamento de imprecisões para modelagem consistente \cite{moutarlier2006experimental}.
  \end{itemize}
\end{frame}

%------------------------------------------------
\section{Importância do Mapeamento Incremental}
%------------------------------------------------
\begin{frame}
  \frametitle{Importância do Mapeamento Incremental}
  \begin{itemize}
      \item Operações-chave: Percepção, reconhecimento, decisão e navegação.
      \item Impacto dos erros sensoriais acumulativos.
      \item Acúmulo de Erros: Como os erros sensoriais se acumulam durante o processo de percepção-modelagem.
      \item Efeito de Erros Anteriores: Impacto nos cálculos atuais
  \end{itemize}
\end{frame}

%------------------------------------------------
\section{Abordagem Teórica}
%------------------------------------------------
\begin{frame}
  \frametitle{Abordagem Teórica}
  \begin{itemize}
      \item Abordagem Rigorosa: Uso de filtragem de Kalman para fusão de dados sensoriais estocásticos.
      \item Desafios: Instabilidade devido à natureza não linear do problema, especialmente sem uma boa calibração da variância de ruído do sensor.
      \item Solução Proposta: Dentro do formalismo apresentado, maneiras de contornar o problema.
  \end{itemize}
\end{frame}
  
%------------------------------------------------
\section{Modelagem Local e Global}
%------------------------------------------------
\begin{frame}
  \frametitle{Modelagem Local e Global}
  \begin{itemize}
      \item Necessidade de Modelagem Local: Conversão de dados brutos do sensor em um modelo local para correspondência com um modelo global.
      \item Desafio com Dados Brutos: Dados brutos de sensores, como um localizador a laser, são geralmente pobres e instáveis para uso direto.
      \item Segmentação e Extração de Características: Processos usados para criar uma representação de nível superior a partir dos dados brutos.
  \end{itemize}
\end{frame}
  

%------------------------------------------------
\section{Ferramentas Matemáticas para Tratar Inexatidões dos Sensores}
%------------------------------------------------
  \begin{frame}
  \frametitle{Ferramentas Matemáticas para Tratar Inexatidões dos Sensores}
  \begin{itemize}
      \item Uso de Sensores: Diferentes tipos de sensores usados para mapeamento incremental.
      \item Problemas com Erros de Movimento: Como os erros de movimento se acumulam com as imprecisões de percepção.
      \item Abordagem Proposta: Uma abordagem baseada em métodos estocásticos para construir uma representação consistente do ambiente.
  \end{itemize}
  \end{frame}
  

%------------------------------------------------
\section{Representação Básica e Posicionamento do Robô}
%------------------------------------------------
  \begin{frame}
  \frametitle{Representação Básica e Posicionamento do Robô}
  \begin{itemize}
      \item Comparação entre representação relacional e de localização.
      \item Escolha de um quadro de referência único.
  \end{itemize}
  \end{frame}
  

%------------------------------------------------
\section{Avaliação de Inexatidões e Calibração}
%------------------------------------------------
  \begin{frame}
  \frametitle{Avaliação de Inexatidões e Calibração}
  \begin{itemize}
      \item Avaliação da deriva odometria.
      \item Calibração do sensor laser.
  \end{itemize}
  \end{frame}
  

%------------------------------------------------
\section{Estratégias de Correspondência e Atualização do Modelo}
%------------------------------------------------
  \begin{frame}
  \frametitle{Estratégias de Correspondência e Atualização do Modelo}
  \begin{itemize}
      \item Correção de posição e correspondência de segmentos.
      \item Uso de variações para correspondência de segmentos.
  \end{itemize}
  \end{frame}
  

%------------------------------------------------
\section{Considerações finais}
%------------------------------------------------
  \begin{frame}
  \frametitle{Considerações finais}
  \begin{itemize}
      \item Resultados e capacidade do sistema.
      \item Desafios futuros e agradecimentos.
  \end{itemize}
  \end{frame}

%------------------------------------------------
%\section*{Referências}
%------------------------------------------------
\begin{frame}
    % \nocite{*}
    \printbibliography
\end{frame}


\begin{frame}
\titlepage % Print the title page as the first slide
\end{frame}

\end{document}